% (c) 2017 Daniele Zambelli - daniele.zambelli@gmail.com
% 
% I grafici per il capitolo relativo alle disequazioni di secondo grado
%

\newcommand{\eqduesf}[4][0]{% Equazione di secondo grado senza formula.
  % Esempio d'uso:
  % \eqdue*{x^2-4 = 0 \sRarrow x^2 = +4}{-2}{+2}
\(
\hspace*{#1mm} #2
 \begin{array}{l}
  \nearrow\\
  \searrow
 \end{array}
 \begin{array}{l}
  x_1 = #3\\[.8em]
  x_2 = #4
 \end{array}
\)
}

\newcommand{\zeripol}[1]{% Zeri del polinomio.
  % Esempio d'uso:
  % \zeripol{\eqduesf{x^2-4 = 0 \sRarrow x^2 = +4}{-2}{+2}}
  Zeri del polinomio: #1
}

\newcommand{\grafun}[1]{% Funzione con grafico.
  % Esempio d'uso:
  % \grafun{\funzass{x^2-4}{\segniparauz{-2}{+2}}}
  Grafico della funzione: #1
}

\newcommand{\funzass}[2]{% Funzione e grafico.
  % Esempio d'uso:
  % \funzass{x^2-4}{\segniparauz{-2}{+2}}
  \(y = #1\) \quad \raisebox{-4mm}{#2}
}

% \NewDocumentCommand \duesegni{m s m}{
%   % Segno di un polinomio di primo grado.
%   % Se c'è l'asterisco lo zero viene reso con una crocetta.
%   % Sintassi: \duesegni[*]{primosegno}{secondosegno}
%   % Esempio di chiamata:
%   % \disegno{\duesegni*{-}{+}
%   % \disegno{\duesegni{-}{+}
%   % \disegno{\duesegni{-}*{+}
%   \draw [black] 
%     (-2.5, .3) node{$#1$} (+2.5, .3) node{$#3$};
%   {\IfBooleanTF{#2}
%       {\crocetta{0}{.4}}
%       {\cerchietto{0}{.4}}}
% }
% 
% \NewDocumentCommand \duesegni{s m s m}{
%   % Segno di un polinomio di primo grado.
%   % Se c'è l'asterisco lo zero viene reso con una crocetta.
%   % Sintassi: \segni[*]{primosegno}{secondosegno}
%   % Esempio di chiamata:
%   % \disegno{\duesegni*{-}{+}
%   % \disegno{\duesegni{-}{+}
%   % \disegno{\duesegni{-}*{+}
%   \node at (-2, .4) {$#2$};
%   \IfBooleanF{#3}
%     {\IfBooleanTF{#1}
%       {\crocetta{0}{.4}}
%       {\cerchietto{0}{.4}}}
%   \node at (+2, .4) {$#4$};
% }

\newcommand{\funzioneassociata}[2]{% blocco relativo alla funz. associata
  % Sintassi; \funzioneassociata{funzione}{grafico}
  Funzione Associata:\\ 
  \quad #1 \quad \raisebox{-6mm}{#2}
}

\newcommand{\segnofatt}[4]{% Segno di un fattore 
  \def \ea{#1}
  \def \fa{#2}
  \def \inaccessiblebl{#3}
  \def \segno{#4}
 \begin{minipage}{.33\textwidth}
  E.A.:~\ea
 \end{minipage}
 \begin{minipage}{.23\textwidth}
  F.A.:~\fa
 \end{minipage}
 \begin{minipage}{.33\textwidth}
  \begin{inaccessibleblock}[\inaccessiblebl]
  \segno
\end{inaccessibleblock}
 \end{minipage}
}

\newcommand{\grrsegnoprodottoa}{% Segno del prodotto di due funzioni
  \def \grrdimx{10}
  \def \grrdimy{1.2}
  \def \grrnpunti{3}
  \disegno{
    \trespolosegnidis*{\grrdimx}{\grrdimy}
      {x^2-4, -2x+3, f(x)}{-2, +\dfrac{3}{2}, +2}
    \cerchietti{\grrdimx}{\grrdimy}{\grrnpunti}
      {1/1, 3/1, 2/2, 1/3, 2/3, 3/3}
    \segnitrespolo{\grrdimx}{\grrdimy}{\grrnpunti}
      {{+, -, -, +}, {+, +, -, -}, {+, -, +, -}}
  }
}

\newcommand{\solprodottoa}{% Soluzione disequazione prodotto di funzioni
  \disegno{
    \assex{0}{6}{0}
    \inti{0}{1}{3}{-2}{+\frac{3}{2}}{blue}{blue}
    \ray{0}{6}{5}{+2}{blue}
  }
}

\newcommand{\grrsegnoprodottob}{% Segno del prodotto di due funzioni
  \def \grrdimx{10}
  \def \grrdimy{1.2}
  \def \grrnpunti{3}
  \disegno{
    \trespolosegnidis*{\grrdimx}{\grrdimy}
    {3x^2-8x-3, x, f(x)}{-\dfrac{1}{3}, 0, +3}
    \cerchietti{\grrdimx}{\grrdimy}{\grrnpunti}
      {1/1, 3/1, 2/2, 1/3, 2/3, 3/3}
    \segnitrespolo{\grrdimx}{\grrdimy}{\grrnpunti}
      {{+, -, -, +}, {-, -, +, +}, {-, +, -, +}}
  }
}

\newcommand{\solprodottob}{% Soluzione disequazione prodotto di funzioni
  \disegno{
    \assex{0}{6}{0}
    \ray{0}{-0.3}{1.5}{-\frac{1}{3}}{white}
    \inti{0}{3}{5}{0}{+3}{white}{white}
  }
}

\newcommand{\segnofrazionea}{% Segno del prodotto di due funzioni
  \def \grrdimx{10}
  \def \grrdimy{1.2}
  \def \grrnpunti{3}
  \disegno{
    \trespolosegnidis*{\grrdimx}{\grrdimy}
    {-x +29, x^2-3x-28, f(x)}{-4, +7, +29}
    \cerchietti{\grrdimx}{\grrdimy}{\grrnpunti}
      {3/1, 3/3}
    \crocette{\grrdimx}{\grrdimy}{\grrnpunti}
      {1/2, 2/2, 1/3, 2/3}
    \segnitrespolo{\grrdimx}{\grrdimy}{\grrnpunti}
      {{+, +, +, -}, {+, -, +, +}, {+, -, +, -}}
  }
}

\newcommand{\solfrazionea}{% Soluzione disequazione prodotto di funzioni
  \disegno{
    \assex{-4}{+4}{0}
    \ray{0}{-4}{-2}{-4}{white}
    \inti{0}{0}{2}{7}{29}{white}{blue}
  }
}

\newcommand{\segnofrazioneb}{% Segno del quoziente di due funzioni
  \def \grrdimx{12}
  \def \grrdimy{1.2}
  \def \grrnpunti{4}
  \disegno{
    \trespolosegnidis*{\grrdimx}{\grrdimy}
      {-2x^2+x+3, 4x^2-1, f(x)}
      {-1, -\dfrac{1}{2}, +\dfrac{1}{2}, +\dfrac{3}{2}}
    \cerchietti{\grrdimx}{\grrdimy}{\grrnpunti}
      {1/1, 4/1, 1/3, 4/3}
    \crocette{\grrdimx}{\grrdimy}{\grrnpunti}
      {2/2, 3/2, 2/3, 3/3}
    \segnitrespolo{\grrdimx}{\grrdimy}{\grrnpunti}
      {{-, +, +, +, -}, {+, +, -, +, +}, {-, +, -, +, -}}
  }
}

\newcommand{\segnofrazionec}{% Segno del quoziente di due funzioni
  \def \grrdimx{12}
  \def \grrdimy{1.2}
  \def \grrnpunti{4}
  \disegno{
    \trespolosegnidis*{\grrdimx}{\grrdimy}
    {-6x^2 +2x +1, -2x^2 +x +1, f(x)}
    {-\dfrac{1}{2}, \dfrac{1 - \sqrt{7}}{6}, \dfrac{1 + \sqrt{7}}{6}, +1}
    \cerchietti{\grrdimx}{\grrdimy}{\grrnpunti}
      {2/1, 3/1, 2/3, 3/3}
    \crocette{\grrdimx}{\grrdimy}{\grrnpunti}
      {1/2, 4/2, 1/3, 4/3}
    \segnitrespolo{\grrdimx}{\grrdimy}{\grrnpunti}
      {{-, -, +, -, -}, {-, +, +, +, -}, {+, -, +, -, +}}
  }
}

\newcommand{\solfrazionec}{% Soluzione disequazione prodotto di funzioni
  \disegno{
    \assex{-5}{+5}{0}
    \ray{0}{-5}{-3}{-\dfrac{1}{2}}{white}
    \inti{0}{-1}{2}
      {\dfrac{1 - \sqrt{7}}{6}}{\dfrac{1 + \sqrt{7}}{6}}{blue}{blue}
    \ray{0}{5}{4}{+1}{white}
  }
}

\newcommand{\segnosistemaaa}{% Segno prima disequazione del primo sistema
  \def \grrdimx{10}
  \def \grrdimy{1.2}
  \def \grrnpunti{3}
  \disegno{
    \trespolosegnidis*{\grrdimx}{\grrdimy}
    {2x+3, x^2-9, f(x)}{-3, -\frac{3}{2}, +3}
    \cerchietti{\grrdimx}{\grrdimy}{\grrnpunti}
      {2/1, 1/2, 3/2, 1/3, 2/3, 3/3}
    \segnitrespolo{\grrdimx}{\grrdimy}{\grrnpunti}
      {{-, -, +, +}, {+, -, -, +}, {-, +, -, +}}
    \ray{-3.6}{0}{2.5}{}{white}
    \inti{-3.6}{5}{7.5}{}{}{white}{white}
  }
}

\newcommand{\segnosistemaab}{% Segno del prodotto di due funzioni
  \def \grrdimx{12}
  \def \grrdimy{1.2}
  \def \grrnpunti{4}
  \disegno{
    \trespolosegnidis*{\grrdimx}{\grrdimy}
      {-x^2 -2x+15, 4x^2-9, f(x)}{-5, -\frac{3}{2}, +\frac{3}{2}, +3}
    \cerchietti{\grrdimx}{\grrdimy}{\grrnpunti}
      {1/1, 4/1, 1/3, 4/3}
    \crocette{\grrdimx}{\grrdimy}{\grrnpunti}
      {2/2, 3/2, 2/3, 3/3}
    \segnitrespolo{\grrdimx}{\grrdimy}{\grrnpunti}
      {{-, +, +, +, -}, {+, +, -, +, +}, {-, +, -, +, -}}
    \ray{-3.6}{0}{2.4}{}{blue}
    \inti{-3.6}{4.8}{7.2}{}{}{white}{white}
    \ray{-3.6}{12}{9.6}{}{blue}
  }
}

\newcommand{\grrsistemaa}{% Intersezione soluzioni sistema a
  \def \grrdimx{12}
  \def \grrdimy{1.2}
  \def \grrnpunti{5}
  \disegno{
    \trespolosegnidis{\grrdimx}{\grrdimy}
      {D_1, D_2, S}{-5, -3, -\frac{3}{2}, +\frac{3}{2}, +3}
    \ray{-1.2}{-0.3}{4}{}{white}
    \inti{-1.2}{6}{10}{}{}{white}{white}
    \ray{-2.4}{-0.3}{2}{}{blue}
    \inti{-2.4}{6}{8}{}{}{white}{white}
    \ray{-2.4}{12}{10}{}{blue}
    \ray{-3.6}{-0.3}{2}{}{blue}
    \inti{-3.6}{6}{8}{}{}{white}{white}
  }
}

\newcommand{\segnosistemaba}{% Segno prima disequazione del primo sistema
  \def \grrdimx{6}
  \def \grrdimy{1.2}
  \def \grrnpunti{2}
  \disegno{
    \trespolosegnidis*{\grrdimx}{\grrdimy}
      {-x^2 +3x +21, x-7, f(x)}{-4, +7}
    \cerchietti{\grrdimx}{\grrdimy}{\grrnpunti}
      {1/1, 2/1, 1/3}
    \crocette{\grrdimx}{\grrdimy}{\grrnpunti}
      {2/2, 2/3}
    \segnitrespolo{\grrdimx}{\grrdimy}{\grrnpunti}
      {{-, +, -}, {-, -, +}, {+, -, -}}
    \ray{-3.6}{-0.3}{2}{}{blue}
  }
}

\newcommand{\segnosistemabb}{% Segno del prodotto di due funzioni
  \def \grrdimx{8}
  \def \grrdimy{1.2}
  \def \grrnpunti{3}
  \disegno{
    \trespolosegnidis*{\grrdimx}{\grrdimy}
      {-x -5, -x^2 +16, f(x)}{-5, -4, +4}
    \cerchietti{\grrdimx}{\grrdimy}{\grrnpunti}
      {1/1, 1/3}
    \crocette{\grrdimx}{\grrdimy}{\grrnpunti}
      {2/2, 3/2, 2/3, 3/3}
    \segnitrespolo{\grrdimx}{\grrdimy}{\grrnpunti}
      {{+, -, -, -}, {-, -, +, -}, {-, +, -, +}}
    \inti{-3.6}{2}{4}{}{}{blue}{white}
    \ray{-3.6}{8}{6}{}{white}
  }
}

\newcommand{\grrsistemab}{% Intersezione soluzioni sistema a
  \def \grrdimx{8}
  \def \grrdimy{1.2}
  \def \grrnpunti{3}
  \disegno{
    \trespolosegnidis{\grrdimx}{\grrdimy}
      {C_1, C_2, S}{-5, -4, +4}
    \ray{-1.2}{-0.3}{4}{}{blue}
    \inti{-2.4}{2}{4}{}{}{blue}{white}
    \ray{-2.4}{8}{6}{}{white}
    \inti{-2.4}{2}{4}{}{}{blue}{white}
    \inti{-3.6}{2}{4}{}{}{blue}{white}
  }
}


\newcommand{\parabole}{% Varie parabole nel piano cartesiano
  \disegno[8]{
    \rcom{-5}{+5}{-4}{+4}{gray!50, very thin, step=1}
    \tkzInit[xmin=-5.3,xmax=+5.3,ymin=-4.3,ymax=+4.3]
    \tkzFct[domain=-1.5:1.5,thick,color=darkgray]{2*x*x-1}
    \node [inner sep=0pt, circle, fill=gray!20] (a) at (-1.1, 2.7) {B};
    \tkzFct[domain=0:3,thick,color=blue]{-2*x*x+6*x-4};
    \node [inner sep=0pt, circle, fill=gray!20] (a) at (.6, -3) {D};
    \tkzFct[domain=-3.3:-.7,thick,color=RedOrange]{-2*x*x-8*x-8.5};
    \node [inner sep=0pt, circle, fill=gray!20] (a) at (-3, -3.5) {E};
    \tkzFct[domain=1.5:4.5,thick,color=purple]{2*x*x-12*x+18};
    \node [inner sep=0pt, circle, fill=gray!20] (a) at (2, 3.7) {C};
    \tkzFct[domain=-4.5:-1.5,thick,color=olive]{2*x*x+12*x+19};
    \node [inner sep=0pt, circle, fill=gray!20] (a) at (-3.8, 3.5) {A};
  }
}
